% ============================================================================
% (c) 2025 Mohamed EL AFRIT -- IPSSI
% Licence CC BY-NC-SA 4.0
% https://creativecommons.org/licenses/by-nc-sa/4.0/deed.fr
%
% FICHE D'EXERCICES (enonces) -- Jour 4 : Reseaux de neurones
% ============================================================================
% ============================================================================
% © 2025 Mohamed EL AFRIT — IPSSI
% Ce contenu est distribué sous licence Creative Commons BY-NC-SA 4.0
% https://creativecommons.org/licenses/by-nc-sa/4.0/deed.fr
% ============================================================================
%
% TEMPLATE COMMUN — Mathématiques appliquées à l'Intelligence Artificielle
% Ce fichier est le préambule partagé de TOUS les documents LaTeX du projet.
%
% Utilisation :
%   % ============================================================================
% © 2025 Mohamed EL AFRIT — IPSSI
% Ce contenu est distribué sous licence Creative Commons BY-NC-SA 4.0
% https://creativecommons.org/licenses/by-nc-sa/4.0/deed.fr
% ============================================================================
%
% TEMPLATE COMMUN — Mathématiques appliquées à l'Intelligence Artificielle
% Ce fichier est le préambule partagé de TOUS les documents LaTeX du projet.
%
% Utilisation :
%   % ============================================================================
% © 2025 Mohamed EL AFRIT — IPSSI
% Ce contenu est distribué sous licence Creative Commons BY-NC-SA 4.0
% https://creativecommons.org/licenses/by-nc-sa/4.0/deed.fr
% ============================================================================
%
% TEMPLATE COMMUN — Mathématiques appliquées à l'Intelligence Artificielle
% Ce fichier est le préambule partagé de TOUS les documents LaTeX du projet.
%
% Utilisation :
%   \input{template_ipssi}          % Importer ce préambule
%   \jourNumero{1}                  % Définir le numéro du jour (1, 2, 3 ou 4)
%   \jourTheme{Données et représentation mathématique}
%   \begin{document}
%   \pagegarde{Fiche de synthèse}{1}{Données et représentation mathématique}
%   ...
%   \end{document}
% ============================================================================

% --- Classe de document ---
\documentclass[11pt,a4paper]{article}

% ============================================================================
% ENCODAGE ET LANGUE
% ============================================================================
\usepackage[utf8]{inputenc}
\usepackage[T1]{fontenc}
\usepackage[french]{babel}

% ============================================================================
% MATHÉMATIQUES
% ============================================================================
\usepackage{amsmath,amssymb,amsfonts,mathtools}

% Commandes mathématiques personnalisées (conventions du cours)
\newcommand{\vx}{\mathbf{x}}          % vecteur x
\newcommand{\vw}{\mathbf{w}}          % vecteur w (poids)
\newcommand{\vy}{\mathbf{y}}          % vecteur y
\newcommand{\mX}{\mathbf{X}}          % matrice X
\newcommand{\mW}{\mathbf{W}}          % matrice W
\newcommand{\yhat}{\hat{y}}           % prédiction
\newcommand{\pscal}[2]{\langle #1, #2 \rangle} % produit scalaire
\newcommand{\gradJ}{\nabla J(\boldsymbol{\theta})} % gradient de J

% ============================================================================
% MISE EN PAGE
% ============================================================================
\usepackage[margin=2.5cm]{geometry}
\usepackage{setspace}
\setstretch{1.15}                      % interligne légèrement aéré

% ============================================================================
% COULEURS
% ============================================================================
\usepackage[dvipsnames,svgnames,x11names]{xcolor}

% Palette principale du projet
\definecolor{ipssiBleu}{HTML}{3B82F6}       % Définitions
\definecolor{ipssiVert}{HTML}{22C55E}       % À retenir
\definecolor{ipssiOrange}{HTML}{F97316}     % Attention / Piège
\definecolor{ipssiViolet}{HTML}{A855F7}     % Intuition
\definecolor{ipssiCyan}{HTML}{06B6D4}       % Exemples
\definecolor{ipssiGrisClair}{HTML}{F3F4F6}  % Fond code
\definecolor{ipssiGrisFonce}{HTML}{1F2937}  % Texte code
\definecolor{ipssiFondPage}{HTML}{FAFAFA}   % Fond léger page de garde
\definecolor{ipssiRouge}{HTML}{E74C3C}      % Classe 0 (salaire ≤ 45k€)
\definecolor{ipssiVertData}{HTML}{2ECC71}   % Classe 1 (salaire > 45k€)
\definecolor{ipssiBleuDroite}{HTML}{3498DB} % Droite de régression

% ============================================================================
% ENCADRÉS PÉDAGOGIQUES (tcolorbox)
% ============================================================================
\usepackage[most]{tcolorbox}

% Style de base partagé
\tcbset{
    ipssibase/.style={
        enhanced,
        breakable,
        left=4mm,
        right=4mm,
        top=3mm,
        bottom=3mm,
        boxrule=0pt,
        frame hidden,
        borderline west={3pt}{0pt}{#1},
        colback=#1!5,
        colframe=#1,
        fonttitle=\bfseries\sffamily,
        coltitle=#1!80!black,
        attach boxed title to top left={yshift=-2mm, xshift=4mm},
        boxed title style={
            colback=#1!15,
            colframe=#1,
            boxrule=0.5pt,
            arc=2pt,
        },
        arc=2pt,
        outer arc=2pt,
    }
}

% Icônes TikZ pour les encadrés (compatible pdflatex, pas d'emoji Unicode)
\newcommand{\icoDefinition}{%
    \tikz[baseline=-0.3em]{\draw[ipssiBleu, thick] (0,0) -- (0.3,0) -- (0.3,0.35) -- (0,0.35) -- cycle;
    \draw[ipssiBleu, thick] (0.05,0.1) -- (0.25,0.1); \draw[ipssiBleu, thick] (0.05,0.2) -- (0.2,0.2);}%
}
\newcommand{\icoRetenir}{%
    \tikz[baseline=-0.2em]{\draw[ipssiVert, thick, fill=ipssiVert!20] (0.15,0) -- (0.3,0.3) -- (0,0.3) -- cycle;}%
}
\newcommand{\icoAttention}{%
    \tikz[baseline=-0.2em]{\draw[ipssiOrange, thick] (0.15,0.35) -- (0,0) -- (0.3,0) -- cycle;
    \node[ipssiOrange, font=\tiny\bfseries] at (0.15,0.12) {!};}%
}
\newcommand{\icoIntuition}{%
    \tikz[baseline=-0.2em]{\draw[ipssiViolet, thick, fill=ipssiViolet!15]
    (0.15,0.35) -- (0.08,0.2) -- (0.08,0.15) -- (0.22,0.15) -- (0.22,0.2) -- cycle;
    \draw[ipssiViolet, thick] (0.11,0.12) -- (0.11,0.05); \draw[ipssiViolet, thick] (0.19,0.12) -- (0.19,0.05);}%
}
\newcommand{\icoExemple}{%
    \tikz[baseline=-0.2em]{\draw[ipssiCyan, thick] (0,0) rectangle (0.3,0.3);
    \draw[ipssiCyan, thick] (0,0.15) -- (0.3,0.15); \draw[ipssiCyan, thick] (0.15,0) -- (0.15,0.3);}%
}
\newcommand{\icoPython}{%
    \tikz[baseline=-0.2em]{\node[font=\small\bfseries\ttfamily, ipssiVert!70!black] at (0,0) {\textgreater\textgreater};}%
}

% Définition (bleu)
\newtcolorbox{definitionbox}[1][D\'{e}finition]{
    ipssibase=ipssiBleu,
    title={\icoDefinition\; #1},
}

% À retenir (vert)
\newtcolorbox{retenirbox}[1][\`{A} retenir]{
    ipssibase=ipssiVert,
    title={\icoRetenir\; #1},
}

% Attention / Piège (orange)
\newtcolorbox{attentionbox}[1][Attention]{
    ipssibase=ipssiOrange,
    title={\icoAttention\; #1},
}

% Intuition (violet)
\newtcolorbox{intuitionbox}[1][Intuition]{
    ipssibase=ipssiViolet,
    title={\icoIntuition\; #1},
}

% Exemple (cyan)
\newtcolorbox{exemplebox}[1][Exemple]{
    ipssibase=ipssiCyan,
    title={\icoExemple\; #1},
}

% Code Python (style terminal)
\newtcolorbox{pythonbox}[1][Code Python]{
    enhanced,
    breakable,
    left=4mm, right=4mm, top=3mm, bottom=3mm,
    boxrule=0.5pt,
    colback=ipssiGrisClair,
    colframe=ipssiGrisFonce!40,
    fonttitle=\bfseries\ttfamily\small,
    coltitle=ipssiGrisFonce,
    title={\icoPython\; #1},
    attach boxed title to top left={yshift=-2mm, xshift=4mm},
    boxed title style={colback=ipssiGrisClair, colframe=ipssiGrisFonce!40, boxrule=0.5pt, arc=2pt},
    arc=2pt, outer arc=2pt,
}

% ============================================================================
% CODE PYTHON (listings)
% ============================================================================
\usepackage{listings}

\lstdefinestyle{python}{
    language=Python,
    basicstyle=\ttfamily\small,
    keywordstyle=\color{ipssiBleu}\bfseries,
    stringstyle=\color{ipssiVert!70!black},
    commentstyle=\color{gray}\itshape,
    numberstyle=\tiny\color{gray},
    numbers=left,
    numbersep=8pt,
    backgroundcolor=\color{ipssiGrisClair},
    frame=none,
    rulecolor=\color{ipssiGrisFonce!20},
    breaklines=true,
    breakatwhitespace=true,
    tabsize=4,
    showstringspaces=false,
    captionpos=b,
    aboveskip=8pt,
    belowskip=8pt,
    xleftmargin=10pt,
    framexleftmargin=10pt,
    morekeywords={as, with, True, False, None, self, np, plt, import, from},
    literate=
        {é}{{\'e}}1 {è}{{\`e}}1 {ê}{{\^e}}1 {ë}{{\"e}}1
        {à}{{\`a}}1 {â}{{\^a}}1
        {ù}{{\`u}}1 {û}{{\^u}}1
        {ô}{{\^o}}1
        {î}{{\^i}}1 {ï}{{\"i}}1
        {ç}{{\c{c}}}1
        {É}{{\'E}}1 {È}{{\`E}}1
        {→}{{$\rightarrow$}}1
        {≈}{{$\approx$}}1
        {≤}{{$\leq$}}1
        {≥}{{$\geq$}}1
        {★}{{$\bigstar$}}1
        {ŷ}{{$\hat{y}$}}1
        {·}{{$\cdot$}}1,
}

\lstset{style=python}

% Style pour le code inline
\newcommand{\pyinline}[1]{\lstinline[style=python,basicstyle=\ttfamily\small]{#1}}

% ============================================================================
% GRAPHIQUES
% ============================================================================
\usepackage{tikz}
\usepackage{pgfplots}
\pgfplotsset{compat=1.18}
\usetikzlibrary{arrows.meta, calc, positioning, shapes.geometric, decorations.pathreplacing}

% ============================================================================
% TABLEAUX
% ============================================================================
\usepackage{booktabs}
\usepackage{tabularx}
\usepackage{multirow}

% ============================================================================
% LISTES
% ============================================================================
\usepackage{enumitem}
\setlist[itemize]{itemsep=2pt, parsep=0pt, topsep=4pt}
\setlist[enumerate]{itemsep=2pt, parsep=0pt, topsep=4pt}

% ============================================================================
% HYPERLIENS
% ============================================================================
\usepackage[
    colorlinks=true,
    linkcolor=ipssiBleu!80!black,
    urlcolor=ipssiBleu!80!black,
    citecolor=ipssiVert!80!black,
    bookmarks=true,
    bookmarksopen=true,
    pdfauthor={Mohamed EL AFRIT},
    pdftitle={Mathématiques appliquées à l'IA — IPSSI},
    pdfsubject={Formation Bachelor 2 — IPSSI 2025-2026},
]{hyperref}

% ============================================================================
% IMAGES
% ============================================================================
\usepackage{graphicx}

% ============================================================================
% VARIABLES PARAMÉTRABLES (jour et thème)
% ============================================================================
\newcommand{\leJour}{X}
\newcommand{\leTheme}{Thème}
\newcommand{\jourNumero}[1]{\renewcommand{\leJour}{#1}}
\newcommand{\jourTheme}[1]{\renewcommand{\leTheme}{#1}}

% ============================================================================
% EN-TÊTE ET PIED DE PAGE (fancyhdr)
% ============================================================================
\usepackage{fancyhdr}

\pagestyle{fancy}
\fancyhf{}

% En-tête
\fancyhead[L]{\small\sffamily\textcolor{gray}{IPSSI — Maths appliquées à l'IA}}
\fancyhead[R]{\small\sffamily\textcolor{gray}{Jour \leJour\ — \leTheme}}
\renewcommand{\headrulewidth}{0.4pt}

% Pied de page
\fancyfoot[L]{\footnotesize\textcolor{gray}{Mohamed EL AFRIT\enspace\textbullet\enspace\url{www.mohamedelafrit.com}\enspace\textbullet\enspace CC BY-NC-SA 4.0}}
\fancyfoot[C]{\footnotesize\textcolor{gray}{\thepage}}
\fancyfoot[R]{\footnotesize\textcolor{gray}{2025–2026}}
\renewcommand{\footrulewidth}{0.2pt}

% Style pour la page de garde (pas d'en-tête)
\fancypagestyle{pagegarde}{
    \fancyhf{}
    \renewcommand{\headrulewidth}{0pt}
    \fancyfoot[C]{\footnotesize\textcolor{gray}{\thepage}}
    \renewcommand{\footrulewidth}{0pt}
}

% ============================================================================
% PAGE DE GARDE — \pagegarde{titre}{jour}{theme}
% ============================================================================
\newcommand{\pagegarde}[3]{%
    \thispagestyle{pagegarde}
    \begin{center}
        \vspace*{1.5cm}

        % Logo / Bandeau institution
        {\Large\sffamily\bfseries\textcolor{ipssiBleu}{IPSSI}}\par
        \vspace{2pt}
        {\normalsize\sffamily 2\textsuperscript{e} année Bachelor Informatique}\par

        \vspace{0.8cm}
        {\color{ipssiBleu}\rule{0.6\textwidth}{1.5pt}}\par
        \vspace{0.8cm}

        % Titre du module
        {\LARGE\sffamily\bfseries Mathématiques appliquées\\[4pt] à l'Intelligence Artificielle}\par

        \vspace{1cm}

        % Titre du document
        \begin{tcolorbox}[
            enhanced,
            colback=ipssiBleu!5,
            colframe=ipssiBleu!50,
            boxrule=1pt,
            arc=4pt,
            width=0.85\textwidth,
            halign=center,
            top=10pt, bottom=10pt,
        ]
            {\huge\sffamily\bfseries\textcolor{ipssiBleu!80!black}{#1}}\par
            \vspace{6pt}
            {\Large\sffamily Jour #2 — #3}
        \end{tcolorbox}

        \vspace{1.5cm}

        % Informations auteur
        {\large\sffamily
            \textbf{Auteur :} Mohamed EL AFRIT\\[4pt]
            \url{www.mohamedelafrit.com}\\[8pt]
            \textbf{Année académique :} 2025–2026
        }\par

        \vspace{1.5cm}

        % Licence Creative Commons
        {\color{ipssiBleu}\rule{0.4\textwidth}{0.5pt}}\par
        \vspace{8pt}
        {\small\sffamily
            \textcopyright\ 2025 Mohamed EL AFRIT — IPSSI\\[2pt]
            Ce contenu est distribué sous licence\\[2pt]
            \textbf{Creative Commons BY-NC-SA 4.0}\\[2pt]
            \url{https://creativecommons.org/licenses/by-nc-sa/4.0/deed.fr}
        }\par
    \end{center}
    \newpage
}

% ============================================================================
% NIVEAUX D'EXERCICES
% ============================================================================

% Étoile compatible pdflatex
\newcommand{\starfull}[1]{\textcolor{#1}{$\bigstar$}}
\newcommand{\starempty}{\textcolor{gray!40}{$\bigstar$}}

% Fondamental (1 étoile)
\newcommand{\exofondamental}{%
    \starfull{ipssiVert}\starempty\starempty\enspace%
    {\small\sffamily\textcolor{ipssiVert}{\textbf{Fondamental}}}%
}

% Intermédiaire (2 étoiles)
\newcommand{\exointermediaire}{%
    \starfull{ipssiOrange}\starfull{ipssiOrange}\starempty\enspace%
    {\small\sffamily\textcolor{ipssiOrange}{\textbf{Interm\'{e}diaire}}}%
}

% Avancé (3 étoiles)
\newcommand{\exoavance}{%
    \starfull{ipssiRouge}\starfull{ipssiRouge}\starfull{ipssiRouge}\enspace%
    {\small\sffamily\textcolor{ipssiRouge}{\textbf{Avanc\'{e}}}}%
}

% Commande d'exercice numéroté
\newcounter{exocount}[section]
\newcommand{\exercice}[1]{%
    \stepcounter{exocount}%
    \vspace{8pt}%
    \noindent\textbf{\sffamily Exercice \theexocount}\enspace — \enspace #1\par
    \vspace{4pt}%
}

% ============================================================================
% COMMANDES UTILITAIRES
% ============================================================================

% Séparateur visuel entre sections
\newcommand{\separateur}{%
    \vspace{6pt}
    \begin{center}
        {\color{ipssiBleu!30}\rule{0.3\textwidth}{0.5pt}}
    \end{center}
    \vspace{6pt}
}

% Mot-clé en gras coloré
\newcommand{\motcle}[1]{\textbf{\textcolor{ipssiBleu!80!black}{#1}}}

% Formule encadrée importante
\newcommand{\formule}[1]{%
    \begin{center}
        \fcolorbox{ipssiBleu!50}{ipssiBleu!5}{%
            \quad$\displaystyle #1$\quad%
        }
    \end{center}
}

% ============================================================================
% FIN DU PRÉAMBULE
% ============================================================================
          % Importer ce préambule
%   \jourNumero{1}                  % Définir le numéro du jour (1, 2, 3 ou 4)
%   \jourTheme{Données et représentation mathématique}
%   \begin{document}
%   \pagegarde{Fiche de synthèse}{1}{Données et représentation mathématique}
%   ...
%   \end{document}
% ============================================================================

% --- Classe de document ---
\documentclass[11pt,a4paper]{article}

% ============================================================================
% ENCODAGE ET LANGUE
% ============================================================================
\usepackage[utf8]{inputenc}
\usepackage[T1]{fontenc}
\usepackage[french]{babel}

% ============================================================================
% MATHÉMATIQUES
% ============================================================================
\usepackage{amsmath,amssymb,amsfonts,mathtools}

% Commandes mathématiques personnalisées (conventions du cours)
\newcommand{\vx}{\mathbf{x}}          % vecteur x
\newcommand{\vw}{\mathbf{w}}          % vecteur w (poids)
\newcommand{\vy}{\mathbf{y}}          % vecteur y
\newcommand{\mX}{\mathbf{X}}          % matrice X
\newcommand{\mW}{\mathbf{W}}          % matrice W
\newcommand{\yhat}{\hat{y}}           % prédiction
\newcommand{\pscal}[2]{\langle #1, #2 \rangle} % produit scalaire
\newcommand{\gradJ}{\nabla J(\boldsymbol{\theta})} % gradient de J

% ============================================================================
% MISE EN PAGE
% ============================================================================
\usepackage[margin=2.5cm]{geometry}
\usepackage{setspace}
\setstretch{1.15}                      % interligne légèrement aéré

% ============================================================================
% COULEURS
% ============================================================================
\usepackage[dvipsnames,svgnames,x11names]{xcolor}

% Palette principale du projet
\definecolor{ipssiBleu}{HTML}{3B82F6}       % Définitions
\definecolor{ipssiVert}{HTML}{22C55E}       % À retenir
\definecolor{ipssiOrange}{HTML}{F97316}     % Attention / Piège
\definecolor{ipssiViolet}{HTML}{A855F7}     % Intuition
\definecolor{ipssiCyan}{HTML}{06B6D4}       % Exemples
\definecolor{ipssiGrisClair}{HTML}{F3F4F6}  % Fond code
\definecolor{ipssiGrisFonce}{HTML}{1F2937}  % Texte code
\definecolor{ipssiFondPage}{HTML}{FAFAFA}   % Fond léger page de garde
\definecolor{ipssiRouge}{HTML}{E74C3C}      % Classe 0 (salaire ≤ 45k€)
\definecolor{ipssiVertData}{HTML}{2ECC71}   % Classe 1 (salaire > 45k€)
\definecolor{ipssiBleuDroite}{HTML}{3498DB} % Droite de régression

% ============================================================================
% ENCADRÉS PÉDAGOGIQUES (tcolorbox)
% ============================================================================
\usepackage[most]{tcolorbox}

% Style de base partagé
\tcbset{
    ipssibase/.style={
        enhanced,
        breakable,
        left=4mm,
        right=4mm,
        top=3mm,
        bottom=3mm,
        boxrule=0pt,
        frame hidden,
        borderline west={3pt}{0pt}{#1},
        colback=#1!5,
        colframe=#1,
        fonttitle=\bfseries\sffamily,
        coltitle=#1!80!black,
        attach boxed title to top left={yshift=-2mm, xshift=4mm},
        boxed title style={
            colback=#1!15,
            colframe=#1,
            boxrule=0.5pt,
            arc=2pt,
        },
        arc=2pt,
        outer arc=2pt,
    }
}

% Icônes TikZ pour les encadrés (compatible pdflatex, pas d'emoji Unicode)
\newcommand{\icoDefinition}{%
    \tikz[baseline=-0.3em]{\draw[ipssiBleu, thick] (0,0) -- (0.3,0) -- (0.3,0.35) -- (0,0.35) -- cycle;
    \draw[ipssiBleu, thick] (0.05,0.1) -- (0.25,0.1); \draw[ipssiBleu, thick] (0.05,0.2) -- (0.2,0.2);}%
}
\newcommand{\icoRetenir}{%
    \tikz[baseline=-0.2em]{\draw[ipssiVert, thick, fill=ipssiVert!20] (0.15,0) -- (0.3,0.3) -- (0,0.3) -- cycle;}%
}
\newcommand{\icoAttention}{%
    \tikz[baseline=-0.2em]{\draw[ipssiOrange, thick] (0.15,0.35) -- (0,0) -- (0.3,0) -- cycle;
    \node[ipssiOrange, font=\tiny\bfseries] at (0.15,0.12) {!};}%
}
\newcommand{\icoIntuition}{%
    \tikz[baseline=-0.2em]{\draw[ipssiViolet, thick, fill=ipssiViolet!15]
    (0.15,0.35) -- (0.08,0.2) -- (0.08,0.15) -- (0.22,0.15) -- (0.22,0.2) -- cycle;
    \draw[ipssiViolet, thick] (0.11,0.12) -- (0.11,0.05); \draw[ipssiViolet, thick] (0.19,0.12) -- (0.19,0.05);}%
}
\newcommand{\icoExemple}{%
    \tikz[baseline=-0.2em]{\draw[ipssiCyan, thick] (0,0) rectangle (0.3,0.3);
    \draw[ipssiCyan, thick] (0,0.15) -- (0.3,0.15); \draw[ipssiCyan, thick] (0.15,0) -- (0.15,0.3);}%
}
\newcommand{\icoPython}{%
    \tikz[baseline=-0.2em]{\node[font=\small\bfseries\ttfamily, ipssiVert!70!black] at (0,0) {\textgreater\textgreater};}%
}

% Définition (bleu)
\newtcolorbox{definitionbox}[1][D\'{e}finition]{
    ipssibase=ipssiBleu,
    title={\icoDefinition\; #1},
}

% À retenir (vert)
\newtcolorbox{retenirbox}[1][\`{A} retenir]{
    ipssibase=ipssiVert,
    title={\icoRetenir\; #1},
}

% Attention / Piège (orange)
\newtcolorbox{attentionbox}[1][Attention]{
    ipssibase=ipssiOrange,
    title={\icoAttention\; #1},
}

% Intuition (violet)
\newtcolorbox{intuitionbox}[1][Intuition]{
    ipssibase=ipssiViolet,
    title={\icoIntuition\; #1},
}

% Exemple (cyan)
\newtcolorbox{exemplebox}[1][Exemple]{
    ipssibase=ipssiCyan,
    title={\icoExemple\; #1},
}

% Code Python (style terminal)
\newtcolorbox{pythonbox}[1][Code Python]{
    enhanced,
    breakable,
    left=4mm, right=4mm, top=3mm, bottom=3mm,
    boxrule=0.5pt,
    colback=ipssiGrisClair,
    colframe=ipssiGrisFonce!40,
    fonttitle=\bfseries\ttfamily\small,
    coltitle=ipssiGrisFonce,
    title={\icoPython\; #1},
    attach boxed title to top left={yshift=-2mm, xshift=4mm},
    boxed title style={colback=ipssiGrisClair, colframe=ipssiGrisFonce!40, boxrule=0.5pt, arc=2pt},
    arc=2pt, outer arc=2pt,
}

% ============================================================================
% CODE PYTHON (listings)
% ============================================================================
\usepackage{listings}

\lstdefinestyle{python}{
    language=Python,
    basicstyle=\ttfamily\small,
    keywordstyle=\color{ipssiBleu}\bfseries,
    stringstyle=\color{ipssiVert!70!black},
    commentstyle=\color{gray}\itshape,
    numberstyle=\tiny\color{gray},
    numbers=left,
    numbersep=8pt,
    backgroundcolor=\color{ipssiGrisClair},
    frame=none,
    rulecolor=\color{ipssiGrisFonce!20},
    breaklines=true,
    breakatwhitespace=true,
    tabsize=4,
    showstringspaces=false,
    captionpos=b,
    aboveskip=8pt,
    belowskip=8pt,
    xleftmargin=10pt,
    framexleftmargin=10pt,
    morekeywords={as, with, True, False, None, self, np, plt, import, from},
    literate=
        {é}{{\'e}}1 {è}{{\`e}}1 {ê}{{\^e}}1 {ë}{{\"e}}1
        {à}{{\`a}}1 {â}{{\^a}}1
        {ù}{{\`u}}1 {û}{{\^u}}1
        {ô}{{\^o}}1
        {î}{{\^i}}1 {ï}{{\"i}}1
        {ç}{{\c{c}}}1
        {É}{{\'E}}1 {È}{{\`E}}1
        {→}{{$\rightarrow$}}1
        {≈}{{$\approx$}}1
        {≤}{{$\leq$}}1
        {≥}{{$\geq$}}1
        {★}{{$\bigstar$}}1
        {ŷ}{{$\hat{y}$}}1
        {·}{{$\cdot$}}1,
}

\lstset{style=python}

% Style pour le code inline
\newcommand{\pyinline}[1]{\lstinline[style=python,basicstyle=\ttfamily\small]{#1}}

% ============================================================================
% GRAPHIQUES
% ============================================================================
\usepackage{tikz}
\usepackage{pgfplots}
\pgfplotsset{compat=1.18}
\usetikzlibrary{arrows.meta, calc, positioning, shapes.geometric, decorations.pathreplacing}

% ============================================================================
% TABLEAUX
% ============================================================================
\usepackage{booktabs}
\usepackage{tabularx}
\usepackage{multirow}

% ============================================================================
% LISTES
% ============================================================================
\usepackage{enumitem}
\setlist[itemize]{itemsep=2pt, parsep=0pt, topsep=4pt}
\setlist[enumerate]{itemsep=2pt, parsep=0pt, topsep=4pt}

% ============================================================================
% HYPERLIENS
% ============================================================================
\usepackage[
    colorlinks=true,
    linkcolor=ipssiBleu!80!black,
    urlcolor=ipssiBleu!80!black,
    citecolor=ipssiVert!80!black,
    bookmarks=true,
    bookmarksopen=true,
    pdfauthor={Mohamed EL AFRIT},
    pdftitle={Mathématiques appliquées à l'IA — IPSSI},
    pdfsubject={Formation Bachelor 2 — IPSSI 2025-2026},
]{hyperref}

% ============================================================================
% IMAGES
% ============================================================================
\usepackage{graphicx}

% ============================================================================
% VARIABLES PARAMÉTRABLES (jour et thème)
% ============================================================================
\newcommand{\leJour}{X}
\newcommand{\leTheme}{Thème}
\newcommand{\jourNumero}[1]{\renewcommand{\leJour}{#1}}
\newcommand{\jourTheme}[1]{\renewcommand{\leTheme}{#1}}

% ============================================================================
% EN-TÊTE ET PIED DE PAGE (fancyhdr)
% ============================================================================
\usepackage{fancyhdr}

\pagestyle{fancy}
\fancyhf{}

% En-tête
\fancyhead[L]{\small\sffamily\textcolor{gray}{IPSSI — Maths appliquées à l'IA}}
\fancyhead[R]{\small\sffamily\textcolor{gray}{Jour \leJour\ — \leTheme}}
\renewcommand{\headrulewidth}{0.4pt}

% Pied de page
\fancyfoot[L]{\footnotesize\textcolor{gray}{Mohamed EL AFRIT\enspace\textbullet\enspace\url{www.mohamedelafrit.com}\enspace\textbullet\enspace CC BY-NC-SA 4.0}}
\fancyfoot[C]{\footnotesize\textcolor{gray}{\thepage}}
\fancyfoot[R]{\footnotesize\textcolor{gray}{2025–2026}}
\renewcommand{\footrulewidth}{0.2pt}

% Style pour la page de garde (pas d'en-tête)
\fancypagestyle{pagegarde}{
    \fancyhf{}
    \renewcommand{\headrulewidth}{0pt}
    \fancyfoot[C]{\footnotesize\textcolor{gray}{\thepage}}
    \renewcommand{\footrulewidth}{0pt}
}

% ============================================================================
% PAGE DE GARDE — \pagegarde{titre}{jour}{theme}
% ============================================================================
\newcommand{\pagegarde}[3]{%
    \thispagestyle{pagegarde}
    \begin{center}
        \vspace*{1.5cm}

        % Logo / Bandeau institution
        {\Large\sffamily\bfseries\textcolor{ipssiBleu}{IPSSI}}\par
        \vspace{2pt}
        {\normalsize\sffamily 2\textsuperscript{e} année Bachelor Informatique}\par

        \vspace{0.8cm}
        {\color{ipssiBleu}\rule{0.6\textwidth}{1.5pt}}\par
        \vspace{0.8cm}

        % Titre du module
        {\LARGE\sffamily\bfseries Mathématiques appliquées\\[4pt] à l'Intelligence Artificielle}\par

        \vspace{1cm}

        % Titre du document
        \begin{tcolorbox}[
            enhanced,
            colback=ipssiBleu!5,
            colframe=ipssiBleu!50,
            boxrule=1pt,
            arc=4pt,
            width=0.85\textwidth,
            halign=center,
            top=10pt, bottom=10pt,
        ]
            {\huge\sffamily\bfseries\textcolor{ipssiBleu!80!black}{#1}}\par
            \vspace{6pt}
            {\Large\sffamily Jour #2 — #3}
        \end{tcolorbox}

        \vspace{1.5cm}

        % Informations auteur
        {\large\sffamily
            \textbf{Auteur :} Mohamed EL AFRIT\\[4pt]
            \url{www.mohamedelafrit.com}\\[8pt]
            \textbf{Année académique :} 2025–2026
        }\par

        \vspace{1.5cm}

        % Licence Creative Commons
        {\color{ipssiBleu}\rule{0.4\textwidth}{0.5pt}}\par
        \vspace{8pt}
        {\small\sffamily
            \textcopyright\ 2025 Mohamed EL AFRIT — IPSSI\\[2pt]
            Ce contenu est distribué sous licence\\[2pt]
            \textbf{Creative Commons BY-NC-SA 4.0}\\[2pt]
            \url{https://creativecommons.org/licenses/by-nc-sa/4.0/deed.fr}
        }\par
    \end{center}
    \newpage
}

% ============================================================================
% NIVEAUX D'EXERCICES
% ============================================================================

% Étoile compatible pdflatex
\newcommand{\starfull}[1]{\textcolor{#1}{$\bigstar$}}
\newcommand{\starempty}{\textcolor{gray!40}{$\bigstar$}}

% Fondamental (1 étoile)
\newcommand{\exofondamental}{%
    \starfull{ipssiVert}\starempty\starempty\enspace%
    {\small\sffamily\textcolor{ipssiVert}{\textbf{Fondamental}}}%
}

% Intermédiaire (2 étoiles)
\newcommand{\exointermediaire}{%
    \starfull{ipssiOrange}\starfull{ipssiOrange}\starempty\enspace%
    {\small\sffamily\textcolor{ipssiOrange}{\textbf{Interm\'{e}diaire}}}%
}

% Avancé (3 étoiles)
\newcommand{\exoavance}{%
    \starfull{ipssiRouge}\starfull{ipssiRouge}\starfull{ipssiRouge}\enspace%
    {\small\sffamily\textcolor{ipssiRouge}{\textbf{Avanc\'{e}}}}%
}

% Commande d'exercice numéroté
\newcounter{exocount}[section]
\newcommand{\exercice}[1]{%
    \stepcounter{exocount}%
    \vspace{8pt}%
    \noindent\textbf{\sffamily Exercice \theexocount}\enspace — \enspace #1\par
    \vspace{4pt}%
}

% ============================================================================
% COMMANDES UTILITAIRES
% ============================================================================

% Séparateur visuel entre sections
\newcommand{\separateur}{%
    \vspace{6pt}
    \begin{center}
        {\color{ipssiBleu!30}\rule{0.3\textwidth}{0.5pt}}
    \end{center}
    \vspace{6pt}
}

% Mot-clé en gras coloré
\newcommand{\motcle}[1]{\textbf{\textcolor{ipssiBleu!80!black}{#1}}}

% Formule encadrée importante
\newcommand{\formule}[1]{%
    \begin{center}
        \fcolorbox{ipssiBleu!50}{ipssiBleu!5}{%
            \quad$\displaystyle #1$\quad%
        }
    \end{center}
}

% ============================================================================
% FIN DU PRÉAMBULE
% ============================================================================
          % Importer ce préambule
%   \jourNumero{1}                  % Définir le numéro du jour (1, 2, 3 ou 4)
%   \jourTheme{Données et représentation mathématique}
%   \begin{document}
%   \pagegarde{Fiche de synthèse}{1}{Données et représentation mathématique}
%   ...
%   \end{document}
% ============================================================================

% --- Classe de document ---
\documentclass[11pt,a4paper]{article}

% ============================================================================
% ENCODAGE ET LANGUE
% ============================================================================
\usepackage[utf8]{inputenc}
\usepackage[T1]{fontenc}
\usepackage[french]{babel}

% ============================================================================
% MATHÉMATIQUES
% ============================================================================
\usepackage{amsmath,amssymb,amsfonts,mathtools}

% Commandes mathématiques personnalisées (conventions du cours)
\newcommand{\vx}{\mathbf{x}}          % vecteur x
\newcommand{\vw}{\mathbf{w}}          % vecteur w (poids)
\newcommand{\vy}{\mathbf{y}}          % vecteur y
\newcommand{\mX}{\mathbf{X}}          % matrice X
\newcommand{\mW}{\mathbf{W}}          % matrice W
\newcommand{\yhat}{\hat{y}}           % prédiction
\newcommand{\pscal}[2]{\langle #1, #2 \rangle} % produit scalaire
\newcommand{\gradJ}{\nabla J(\boldsymbol{\theta})} % gradient de J

% ============================================================================
% MISE EN PAGE
% ============================================================================
\usepackage[margin=2.5cm]{geometry}
\usepackage{setspace}
\setstretch{1.15}                      % interligne légèrement aéré

% ============================================================================
% COULEURS
% ============================================================================
\usepackage[dvipsnames,svgnames,x11names]{xcolor}

% Palette principale du projet
\definecolor{ipssiBleu}{HTML}{3B82F6}       % Définitions
\definecolor{ipssiVert}{HTML}{22C55E}       % À retenir
\definecolor{ipssiOrange}{HTML}{F97316}     % Attention / Piège
\definecolor{ipssiViolet}{HTML}{A855F7}     % Intuition
\definecolor{ipssiCyan}{HTML}{06B6D4}       % Exemples
\definecolor{ipssiGrisClair}{HTML}{F3F4F6}  % Fond code
\definecolor{ipssiGrisFonce}{HTML}{1F2937}  % Texte code
\definecolor{ipssiFondPage}{HTML}{FAFAFA}   % Fond léger page de garde
\definecolor{ipssiRouge}{HTML}{E74C3C}      % Classe 0 (salaire ≤ 45k€)
\definecolor{ipssiVertData}{HTML}{2ECC71}   % Classe 1 (salaire > 45k€)
\definecolor{ipssiBleuDroite}{HTML}{3498DB} % Droite de régression

% ============================================================================
% ENCADRÉS PÉDAGOGIQUES (tcolorbox)
% ============================================================================
\usepackage[most]{tcolorbox}

% Style de base partagé
\tcbset{
    ipssibase/.style={
        enhanced,
        breakable,
        left=4mm,
        right=4mm,
        top=3mm,
        bottom=3mm,
        boxrule=0pt,
        frame hidden,
        borderline west={3pt}{0pt}{#1},
        colback=#1!5,
        colframe=#1,
        fonttitle=\bfseries\sffamily,
        coltitle=#1!80!black,
        attach boxed title to top left={yshift=-2mm, xshift=4mm},
        boxed title style={
            colback=#1!15,
            colframe=#1,
            boxrule=0.5pt,
            arc=2pt,
        },
        arc=2pt,
        outer arc=2pt,
    }
}

% Icônes TikZ pour les encadrés (compatible pdflatex, pas d'emoji Unicode)
\newcommand{\icoDefinition}{%
    \tikz[baseline=-0.3em]{\draw[ipssiBleu, thick] (0,0) -- (0.3,0) -- (0.3,0.35) -- (0,0.35) -- cycle;
    \draw[ipssiBleu, thick] (0.05,0.1) -- (0.25,0.1); \draw[ipssiBleu, thick] (0.05,0.2) -- (0.2,0.2);}%
}
\newcommand{\icoRetenir}{%
    \tikz[baseline=-0.2em]{\draw[ipssiVert, thick, fill=ipssiVert!20] (0.15,0) -- (0.3,0.3) -- (0,0.3) -- cycle;}%
}
\newcommand{\icoAttention}{%
    \tikz[baseline=-0.2em]{\draw[ipssiOrange, thick] (0.15,0.35) -- (0,0) -- (0.3,0) -- cycle;
    \node[ipssiOrange, font=\tiny\bfseries] at (0.15,0.12) {!};}%
}
\newcommand{\icoIntuition}{%
    \tikz[baseline=-0.2em]{\draw[ipssiViolet, thick, fill=ipssiViolet!15]
    (0.15,0.35) -- (0.08,0.2) -- (0.08,0.15) -- (0.22,0.15) -- (0.22,0.2) -- cycle;
    \draw[ipssiViolet, thick] (0.11,0.12) -- (0.11,0.05); \draw[ipssiViolet, thick] (0.19,0.12) -- (0.19,0.05);}%
}
\newcommand{\icoExemple}{%
    \tikz[baseline=-0.2em]{\draw[ipssiCyan, thick] (0,0) rectangle (0.3,0.3);
    \draw[ipssiCyan, thick] (0,0.15) -- (0.3,0.15); \draw[ipssiCyan, thick] (0.15,0) -- (0.15,0.3);}%
}
\newcommand{\icoPython}{%
    \tikz[baseline=-0.2em]{\node[font=\small\bfseries\ttfamily, ipssiVert!70!black] at (0,0) {\textgreater\textgreater};}%
}

% Définition (bleu)
\newtcolorbox{definitionbox}[1][D\'{e}finition]{
    ipssibase=ipssiBleu,
    title={\icoDefinition\; #1},
}

% À retenir (vert)
\newtcolorbox{retenirbox}[1][\`{A} retenir]{
    ipssibase=ipssiVert,
    title={\icoRetenir\; #1},
}

% Attention / Piège (orange)
\newtcolorbox{attentionbox}[1][Attention]{
    ipssibase=ipssiOrange,
    title={\icoAttention\; #1},
}

% Intuition (violet)
\newtcolorbox{intuitionbox}[1][Intuition]{
    ipssibase=ipssiViolet,
    title={\icoIntuition\; #1},
}

% Exemple (cyan)
\newtcolorbox{exemplebox}[1][Exemple]{
    ipssibase=ipssiCyan,
    title={\icoExemple\; #1},
}

% Code Python (style terminal)
\newtcolorbox{pythonbox}[1][Code Python]{
    enhanced,
    breakable,
    left=4mm, right=4mm, top=3mm, bottom=3mm,
    boxrule=0.5pt,
    colback=ipssiGrisClair,
    colframe=ipssiGrisFonce!40,
    fonttitle=\bfseries\ttfamily\small,
    coltitle=ipssiGrisFonce,
    title={\icoPython\; #1},
    attach boxed title to top left={yshift=-2mm, xshift=4mm},
    boxed title style={colback=ipssiGrisClair, colframe=ipssiGrisFonce!40, boxrule=0.5pt, arc=2pt},
    arc=2pt, outer arc=2pt,
}

% ============================================================================
% CODE PYTHON (listings)
% ============================================================================
\usepackage{listings}

\lstdefinestyle{python}{
    language=Python,
    basicstyle=\ttfamily\small,
    keywordstyle=\color{ipssiBleu}\bfseries,
    stringstyle=\color{ipssiVert!70!black},
    commentstyle=\color{gray}\itshape,
    numberstyle=\tiny\color{gray},
    numbers=left,
    numbersep=8pt,
    backgroundcolor=\color{ipssiGrisClair},
    frame=none,
    rulecolor=\color{ipssiGrisFonce!20},
    breaklines=true,
    breakatwhitespace=true,
    tabsize=4,
    showstringspaces=false,
    captionpos=b,
    aboveskip=8pt,
    belowskip=8pt,
    xleftmargin=10pt,
    framexleftmargin=10pt,
    morekeywords={as, with, True, False, None, self, np, plt, import, from},
    literate=
        {é}{{\'e}}1 {è}{{\`e}}1 {ê}{{\^e}}1 {ë}{{\"e}}1
        {à}{{\`a}}1 {â}{{\^a}}1
        {ù}{{\`u}}1 {û}{{\^u}}1
        {ô}{{\^o}}1
        {î}{{\^i}}1 {ï}{{\"i}}1
        {ç}{{\c{c}}}1
        {É}{{\'E}}1 {È}{{\`E}}1
        {→}{{$\rightarrow$}}1
        {≈}{{$\approx$}}1
        {≤}{{$\leq$}}1
        {≥}{{$\geq$}}1
        {★}{{$\bigstar$}}1
        {ŷ}{{$\hat{y}$}}1
        {·}{{$\cdot$}}1,
}

\lstset{style=python}

% Style pour le code inline
\newcommand{\pyinline}[1]{\lstinline[style=python,basicstyle=\ttfamily\small]{#1}}

% ============================================================================
% GRAPHIQUES
% ============================================================================
\usepackage{tikz}
\usepackage{pgfplots}
\pgfplotsset{compat=1.18}
\usetikzlibrary{arrows.meta, calc, positioning, shapes.geometric, decorations.pathreplacing}

% ============================================================================
% TABLEAUX
% ============================================================================
\usepackage{booktabs}
\usepackage{tabularx}
\usepackage{multirow}

% ============================================================================
% LISTES
% ============================================================================
\usepackage{enumitem}
\setlist[itemize]{itemsep=2pt, parsep=0pt, topsep=4pt}
\setlist[enumerate]{itemsep=2pt, parsep=0pt, topsep=4pt}

% ============================================================================
% HYPERLIENS
% ============================================================================
\usepackage[
    colorlinks=true,
    linkcolor=ipssiBleu!80!black,
    urlcolor=ipssiBleu!80!black,
    citecolor=ipssiVert!80!black,
    bookmarks=true,
    bookmarksopen=true,
    pdfauthor={Mohamed EL AFRIT},
    pdftitle={Mathématiques appliquées à l'IA — IPSSI},
    pdfsubject={Formation Bachelor 2 — IPSSI 2025-2026},
]{hyperref}

% ============================================================================
% IMAGES
% ============================================================================
\usepackage{graphicx}

% ============================================================================
% VARIABLES PARAMÉTRABLES (jour et thème)
% ============================================================================
\newcommand{\leJour}{X}
\newcommand{\leTheme}{Thème}
\newcommand{\jourNumero}[1]{\renewcommand{\leJour}{#1}}
\newcommand{\jourTheme}[1]{\renewcommand{\leTheme}{#1}}

% ============================================================================
% EN-TÊTE ET PIED DE PAGE (fancyhdr)
% ============================================================================
\usepackage{fancyhdr}

\pagestyle{fancy}
\fancyhf{}

% En-tête
\fancyhead[L]{\small\sffamily\textcolor{gray}{IPSSI — Maths appliquées à l'IA}}
\fancyhead[R]{\small\sffamily\textcolor{gray}{Jour \leJour\ — \leTheme}}
\renewcommand{\headrulewidth}{0.4pt}

% Pied de page
\fancyfoot[L]{\footnotesize\textcolor{gray}{Mohamed EL AFRIT\enspace\textbullet\enspace\url{www.mohamedelafrit.com}\enspace\textbullet\enspace CC BY-NC-SA 4.0}}
\fancyfoot[C]{\footnotesize\textcolor{gray}{\thepage}}
\fancyfoot[R]{\footnotesize\textcolor{gray}{2025–2026}}
\renewcommand{\footrulewidth}{0.2pt}

% Style pour la page de garde (pas d'en-tête)
\fancypagestyle{pagegarde}{
    \fancyhf{}
    \renewcommand{\headrulewidth}{0pt}
    \fancyfoot[C]{\footnotesize\textcolor{gray}{\thepage}}
    \renewcommand{\footrulewidth}{0pt}
}

% ============================================================================
% PAGE DE GARDE — \pagegarde{titre}{jour}{theme}
% ============================================================================
\newcommand{\pagegarde}[3]{%
    \thispagestyle{pagegarde}
    \begin{center}
        \vspace*{1.5cm}

        % Logo / Bandeau institution
        {\Large\sffamily\bfseries\textcolor{ipssiBleu}{IPSSI}}\par
        \vspace{2pt}
        {\normalsize\sffamily 2\textsuperscript{e} année Bachelor Informatique}\par

        \vspace{0.8cm}
        {\color{ipssiBleu}\rule{0.6\textwidth}{1.5pt}}\par
        \vspace{0.8cm}

        % Titre du module
        {\LARGE\sffamily\bfseries Mathématiques appliquées\\[4pt] à l'Intelligence Artificielle}\par

        \vspace{1cm}

        % Titre du document
        \begin{tcolorbox}[
            enhanced,
            colback=ipssiBleu!5,
            colframe=ipssiBleu!50,
            boxrule=1pt,
            arc=4pt,
            width=0.85\textwidth,
            halign=center,
            top=10pt, bottom=10pt,
        ]
            {\huge\sffamily\bfseries\textcolor{ipssiBleu!80!black}{#1}}\par
            \vspace{6pt}
            {\Large\sffamily Jour #2 — #3}
        \end{tcolorbox}

        \vspace{1.5cm}

        % Informations auteur
        {\large\sffamily
            \textbf{Auteur :} Mohamed EL AFRIT\\[4pt]
            \url{www.mohamedelafrit.com}\\[8pt]
            \textbf{Année académique :} 2025–2026
        }\par

        \vspace{1.5cm}

        % Licence Creative Commons
        {\color{ipssiBleu}\rule{0.4\textwidth}{0.5pt}}\par
        \vspace{8pt}
        {\small\sffamily
            \textcopyright\ 2025 Mohamed EL AFRIT — IPSSI\\[2pt]
            Ce contenu est distribué sous licence\\[2pt]
            \textbf{Creative Commons BY-NC-SA 4.0}\\[2pt]
            \url{https://creativecommons.org/licenses/by-nc-sa/4.0/deed.fr}
        }\par
    \end{center}
    \newpage
}

% ============================================================================
% NIVEAUX D'EXERCICES
% ============================================================================

% Étoile compatible pdflatex
\newcommand{\starfull}[1]{\textcolor{#1}{$\bigstar$}}
\newcommand{\starempty}{\textcolor{gray!40}{$\bigstar$}}

% Fondamental (1 étoile)
\newcommand{\exofondamental}{%
    \starfull{ipssiVert}\starempty\starempty\enspace%
    {\small\sffamily\textcolor{ipssiVert}{\textbf{Fondamental}}}%
}

% Intermédiaire (2 étoiles)
\newcommand{\exointermediaire}{%
    \starfull{ipssiOrange}\starfull{ipssiOrange}\starempty\enspace%
    {\small\sffamily\textcolor{ipssiOrange}{\textbf{Interm\'{e}diaire}}}%
}

% Avancé (3 étoiles)
\newcommand{\exoavance}{%
    \starfull{ipssiRouge}\starfull{ipssiRouge}\starfull{ipssiRouge}\enspace%
    {\small\sffamily\textcolor{ipssiRouge}{\textbf{Avanc\'{e}}}}%
}

% Commande d'exercice numéroté
\newcounter{exocount}[section]
\newcommand{\exercice}[1]{%
    \stepcounter{exocount}%
    \vspace{8pt}%
    \noindent\textbf{\sffamily Exercice \theexocount}\enspace — \enspace #1\par
    \vspace{4pt}%
}

% ============================================================================
% COMMANDES UTILITAIRES
% ============================================================================

% Séparateur visuel entre sections
\newcommand{\separateur}{%
    \vspace{6pt}
    \begin{center}
        {\color{ipssiBleu!30}\rule{0.3\textwidth}{0.5pt}}
    \end{center}
    \vspace{6pt}
}

% Mot-clé en gras coloré
\newcommand{\motcle}[1]{\textbf{\textcolor{ipssiBleu!80!black}{#1}}}

% Formule encadrée importante
\newcommand{\formule}[1]{%
    \begin{center}
        \fcolorbox{ipssiBleu!50}{ipssiBleu!5}{%
            \quad$\displaystyle #1$\quad%
        }
    \end{center}
}

% ============================================================================
% FIN DU PRÉAMBULE
% ============================================================================


\jourNumero{4}
\jourTheme{R\'{e}seaux de neurones}

\begin{document}

\pagegarde{Fiche d'exercices}{4}{R\'{e}seaux de neurones}

\setstretch{1.08}
\setlength{\parskip}{3pt}
\setlength{\abovedisplayskip}{6pt}
\setlength{\belowdisplayskip}{6pt}

% ======================================================================
% RAPPELS
% ======================================================================
\begin{tcolorbox}[
    enhanced, colback=ipssiCyan!5, colframe=ipssiCyan!50,
    boxrule=0.8pt, arc=4pt,
    title={\sffamily\bfseries Rappels --- Formules du Jour 4},
    coltitle=white, colbacktitle=ipssiCyan!70,
]
\small
\textbf{Neurone :}
$a = \sigma(\vw^T\vx + b)$
\hfill
\textbf{Sigmo\"{\i}de :}
$\sigma(z) = \frac{1}{1+e^{-z}}$
\hfill
\textbf{ReLU :}
$\text{ReLU}(z) = \max(0, z)$

\medskip
\textbf{Forward pass :}
$\mathbf{h} = \sigma(\mathbf{W}^{(1)}\vx + \mathbf{b}^{(1)})$
\;;\;
$\yhat = \sigma(\mathbf{w}^{(2)T}\mathbf{h} + b^{(2)})$

\medskip
\textbf{Mise \`{a} jour :}
$\boldsymbol{\theta} \leftarrow \boldsymbol{\theta} - \alpha \nabla L$
\hfill
\textbf{D\'{e}riv\'{e}e sigmo\"{\i}de :} $\sigma'(z) = \sigma(z)(1-\sigma(z))$

\medskip
\textbf{Valeurs utiles :}
$\sigma(0) = 0{,}5$\;;\;
$\sigma(1) \approx 0{,}73$\;;\;
$\sigma(2) \approx 0{,}88$\;;\;
$\sigma(-1) \approx 0{,}27$\;;\;
$\sigma(-2) \approx 0{,}12$
\end{tcolorbox}

% ======================================================================
\section*{Exercices crayon-papier --- Niveau fondamental}
% ======================================================================

% ----------------------------------------------------------------------
\exercice{\exofondamental\quad --- Sortie d'un neurone}

Un neurone a les poids $\vw = \begin{pmatrix} 0{,}3 \\ 0{,}5 \end{pmatrix}$,
le biais $b = -2$ et utilise l'activation \textbf{sigmo\"{\i}de}.

L'entr\'{e}e est $\vx = \begin{pmatrix} 5 \\ 4 \end{pmatrix}$
(exp\'{e}rience = 5, langages = 4).

\begin{enumerate}[label=\alph*)]
    \item Calculez la somme pond\'{e}r\'{e}e $z = \vw^T\vx + b$.
    \item Calculez l'activation $a = \sigma(z) = \frac{1}{1 + e^{-z}}$.
        \emph{Conseil : calculez d'abord $e^{-z}$ avec une calculatrice ou les valeurs utiles ci-dessus.}
    \item Si le seuil de d\'{e}cision est $0{,}5$, quelle classe le neurone pr\'{e}dit-il ?
    \item On change l'entr\'{e}e : $\vx = \begin{pmatrix} 1 \\ 2 \end{pmatrix}$
        (exp=1, lang=2). Recalculez $z$ et $a$. Quelle classe ?
\end{enumerate}

% ----------------------------------------------------------------------
\exercice{\exofondamental\quad --- Appliquer ReLU et sigmo\"{\i}de}

Calculez les valeurs suivantes en d\'{e}taillant chaque \'{e}tape :

\begin{enumerate}[label=\alph*)]
    \item $\sigma(0)$, $\sigma(2)$, $\sigma(-3)$
    \item $\text{ReLU}(2{,}5)$, $\text{ReLU}(0)$, $\text{ReLU}(-1{,}7)$
    \item $\sigma(\text{ReLU}(-1))$ --- composez les deux fonctions.
    \item $\sigma(\text{ReLU}(3))$ --- m\^{e}me chose.
    \item Pourquoi $\sigma(z)$ ne vaut-il \textbf{jamais} exactement 0 ni exactement 1 ?
\end{enumerate}

% ----------------------------------------------------------------------
\exercice{\exofondamental\quad --- Forward pass d'un r\'{e}seau 2-2-1}

On consid\`{e}re un r\'{e}seau avec 2 entr\'{e}es, 2 neurones cach\'{e}s, 1 sortie.
Tous les neurones utilisent l'activation sigmo\"{\i}de.

\textbf{Poids couche cach\'{e}e :}
\[
    \mathbf{W}^{(1)} = \begin{pmatrix} 0{,}4 & 0{,}6 \\ -0{,}3 & 0{,}5 \end{pmatrix}
    \quad
    \mathbf{b}^{(1)} = \begin{pmatrix} -1 \\ 0{,}5 \end{pmatrix}
\]

\textbf{Poids couche de sortie :}
$\mathbf{w}^{(2)} = \begin{pmatrix} 0{,}7 \\ 0{,}8 \end{pmatrix}$,
$b^{(2)} = -0{,}5$.

Entr\'{e}e : $\vx = \begin{pmatrix} 3 \\ 5 \end{pmatrix}$.

\begin{enumerate}[label=\alph*)]
    \item Calculez les sommes pond\'{e}r\'{e}es de la couche cach\'{e}e :
        $z_1 = \mathbf{W}^{(1)}_{\text{ligne 1}} \cdot \vx + b_1^{(1)}$ et $z_2$.
    \item Appliquez l'activation : $h_1 = \sigma(z_1)$ et $h_2 = \sigma(z_2)$.
    \item Calculez la somme pond\'{e}r\'{e}e de la sortie :
        $z_{\text{out}} = \mathbf{w}^{(2)T}\mathbf{h} + b^{(2)}$.
    \item Calculez la pr\'{e}diction $\yhat = \sigma(z_{\text{out}})$.
    \item Le r\'{e}seau pr\'{e}dit-il classe 0 ou classe 1 ?
\end{enumerate}

% ----------------------------------------------------------------------
\exercice{\exofondamental\quad --- Dimensions des matrices de poids}

Un r\'{e}seau a l'architecture suivante : \textbf{5 entr\'{e}es $\to$ 8 neurones cach\'{e}s $\to$ 3 neurones cach\'{e}s $\to$ 1 sortie}.

\begin{enumerate}[label=\alph*)]
    \item Quelle est la dimension de $\mathbf{W}^{(1)}$ (poids de la 1\textsuperscript{re} couche cach\'{e}e) ?
        Et de $\mathbf{b}^{(1)}$ ?
    \item Quelle est la dimension de $\mathbf{W}^{(2)}$ (poids de la 2\textsuperscript{e} couche) ?
        Et de $\mathbf{b}^{(2)}$ ?
    \item Quelle est la dimension de $\mathbf{w}^{(3)}$ (poids de la sortie) ?
    \item Combien de param\`{e}tres (poids + biais) le r\'{e}seau a-t-il \textbf{au total} ?
    \item Si on augmente la couche cach\'{e}e 1 de 8 \`{a} 16 neurones,
        combien de param\`{e}tres suppl\'{e}mentaires cela ajoute-t-il ?
\end{enumerate}

% ======================================================================
\section*{Exercices crayon-papier --- Niveau interm\'{e}diaire}
% ======================================================================

% ----------------------------------------------------------------------
\exercice{\exointermediaire\quad --- Composer deux fonctions}

On d\'{e}finit deux neurones :
\begin{itemize}
    \item Neurone A : $f_A(x) = \text{ReLU}(2x - 1)$
    \item Neurone B : $g_B(x) = \sigma(3x - 2)$
\end{itemize}

\begin{enumerate}[label=\alph*)]
    \item Calculez $f_A(0)$, $f_A(1)$, $f_A(3)$.
    \item Calculez la composition $g_B(f_A(x))$ pour $x = 0$, $x = 1$, $x = 3$.
        D\'{e}taillez chaque \'{e}tape.
    \item La composition $g_B \circ f_A$ est-elle une fonction lin\'{e}aire ?
        Justifiez.
    \item Que se passerait-il si les deux neurones n'avaient \textbf{pas}
        de fonction d'activation (i.e.\ $f_A(x) = 2x - 1$, $g_B(x) = 3x - 2$) ?
        Calculez la composition. Est-elle lin\'{e}aire ?
\end{enumerate}

% ----------------------------------------------------------------------
\exercice{\exointermediaire\quad --- Pourquoi les activations sont indispensables}

On consid\`{e}re un r\'{e}seau \`{a} 2 couches \textbf{sans} fonction d'activation :
\begin{align*}
    \text{Couche 1 :}\quad & \mathbf{h} = \mathbf{W}^{(1)}\vx + \mathbf{b}^{(1)} \\
    \text{Couche 2 :}\quad & \yhat = \mathbf{w}^{(2)T}\mathbf{h} + b^{(2)}
\end{align*}

\begin{enumerate}[label=\alph*)]
    \item Substituez l'expression de $\mathbf{h}$ dans celle de $\yhat$.
        Simplifiez l'expression.
    \item Montrez que $\yhat$ peut s'\'{e}crire sous la forme $\yhat = \tilde{\vw}^T \vx + \tilde{b}$
        pour un certain $\tilde{\vw}$ et $\tilde{b}$.
    \item Que repr\'{e}sente cette forme ? Le r\'{e}seau \`{a} 2 couches est-il plus
        puissant qu'un simple neurone ?
    \item Conclusion : pourquoi les fonctions d'activation sont-elles
        \textbf{indispensables} dans un r\'{e}seau multicouche ?
\end{enumerate}

% ----------------------------------------------------------------------
\exercice{\exointermediaire\quad --- Analyser un cas de surapprentissage}

Un \'{e}tudiant a entra\^{i}n\'{e} un r\'{e}seau sur le dataset salaire et obtient les
accuracies suivantes au fil des \'{e}poques :

\begin{center}
\renewcommand{\arraystretch}{1.15}
\begin{tabular}{c|c|c}
    \toprule
    \textbf{\'{E}poque} & \textbf{Accuracy entra\^{i}nement} & \textbf{Accuracy validation} \\
    \midrule
    10   & 70\% & 67\% \\
    50   & 85\% & 83\% \\
    200  & 95\% & 87\% \\
    500  & 99\% & 82\% \\
    1000 & 100\% & 75\% \\
    \bottomrule
\end{tabular}
\end{center}

\begin{enumerate}[label=\alph*)]
    \item \`{A} quelle \'{e}poque le mod\`{e}le commence-t-il \`{a} \textbf{surapprentir} ?
        Justifiez.
    \item Quelle \'{e}poque choisiriez-vous comme \textbf{meilleur mod\`{e}le} ? Pourquoi ?
    \item Proposez 3 solutions pour r\'{e}duire le surapprentissage.
    \item Pourquoi une accuracy de 100\% sur l'entra\^{i}nement est-elle
        \textbf{suspecte} plut\^{o}t que rassurante ?
\end{enumerate}

% ======================================================================
\section*{Exercices Python guid\'{e}s --- Niveau interm\'{e}diaire}
% ======================================================================

% ----------------------------------------------------------------------
\exercice{\exointermediaire\quad --- Sigmo\"{\i}de et ReLU en Python \hfill\normalfont\small\textsf{[Python guid\'{e}]}}

\begin{pythonbox}[Squelette --- Fonctions d'activation]
\begin{lstlisting}
import numpy as np
import matplotlib.pyplot as plt

# 1. Implementer la sigmoide
def sigmoid(z):
    return ______  # A COMPLETER

# 2. Implementer ReLU
def relu(z):
    return ______  # A COMPLETER

# 3. Tester sur des valeurs
valeurs = np.array([-3, -1, 0, 1, 2, 5])
print("z       :", valeurs)
print("sigmoid :", np.round(sigmoid(valeurs), 4))
print("relu    :", relu(valeurs))

# 4. Tracer les courbes
z = np.linspace(-6, 6, 200)
fig, (ax1, ax2) = plt.subplots(1, 2, figsize=(12, 4))

ax1.plot(z, ______, 'b-', lw=2)  # A COMPLETER
ax1.set_title("Sigmoide"); ax1.grid(True, alpha=0.3)
ax1.axhline(0.5, color='gray', ls=':')

ax2.plot(z, ______, 'r-', lw=2)  # A COMPLETER
ax2.set_title("ReLU"); ax2.grid(True, alpha=0.3)

plt.tight_layout(); plt.show()

# 5. Derivee de la sigmoide
def sigmoid_deriv(z):
    s = sigmoid(z)
    return ______  # A COMPLETER : s * (1 - s)

print("\nDerivee sigmoid(0) =", sigmoid_deriv(0))
print("Derivee sigmoid(2) =", round(sigmoid_deriv(2), 4))
\end{lstlisting}
\end{pythonbox}

% ----------------------------------------------------------------------
\exercice{\exointermediaire\quad --- Forward pass 2-2-1 en NumPy \hfill\normalfont\small\textsf{[Python guid\'{e}]}}

\begin{pythonbox}[Squelette --- Forward pass]
\begin{lstlisting}
import numpy as np

def sigmoid(z):
    return 1 / (1 + np.exp(-z))

# Entree : un developpeur
x = np.array([5.0, 4.0])

# Poids couche cachee (2 neurones)
W1 = np.array([[0.3, 0.5],
               [-0.2, 0.4]])
b1 = np.array([-2.0, 0.5])

# Poids couche de sortie (1 neurone)
w2 = np.array([0.8, 0.6])
b2 = -0.5

# 1. Couche cachee
z1 = ______           # A COMPLETER : W1 @ x + b1
h  = ______           # A COMPLETER : sigmoid(z1)
print(f"z1 = {z1}")
print(f"h  = {np.round(h, 4)}")

# 2. Couche de sortie
z2 = ______           # A COMPLETER : w2 @ h + b2
y_pred = ______       # A COMPLETER : sigmoid(z2)
print(f"\nz2 = {z2:.4f}")
print(f"y_pred = {y_pred:.4f}")

# 3. Decision
classe = ______       # A COMPLETER : 1 si y_pred > 0.5
print(f"\nClasse predite : {classe}")

# 4. Tester sur tous les developpeurs
print("\n--- Predictions sur 6 developpeurs ---")
devs = np.array([[0,2],[1,1],[5,4],[9,5],[12,7],[18,5]])
for d in devs:
    h_i = sigmoid(W1 @ d + b1)
    y_i = sigmoid(w2 @ h_i + b2)
    print(f"  x={d} -> y_pred={y_i:.3f} "
          f"-> classe {1 if y_i > 0.5 else 0}")
\end{lstlisting}
\end{pythonbox}

% ----------------------------------------------------------------------
\exercice{\exointermediaire\quad --- R\'{e}seau qui apprend XOR \hfill\normalfont\small\textsf{[Python guid\'{e}]}}

\begin{pythonbox}[Squelette --- XOR avec backpropagation]
\begin{lstlisting}
import numpy as np

def sigmoid(z):
    return 1 / (1 + np.exp(-z))
def sigmoid_d(a):
    return a * (1 - a)

# Donnees XOR
X = np.array([[0,0],[0,1],[1,0],[1,1]], dtype=float)
y = np.array([[0],[1],[1],[0]], dtype=float)

# Poids aleatoires (reseau 2-4-1)
np.random.seed(42)
W1 = np.random.randn(2, 4) * 0.5
b1 = np.zeros((1, 4))
W2 = np.random.randn(4, 1) * 0.5
b2 = np.zeros((1, 1))

alpha = 1.0

for epoch in range(5000):
    # Forward
    z1 = ______          # A COMPLETER : X @ W1 + b1
    h  = ______          # A COMPLETER : sigmoid(z1)
    z2 = ______          # A COMPLETER : h @ W2 + b2
    out = ______         # A COMPLETER : sigmoid(z2)

    loss = np.mean((y - out) ** 2)

    # Backward
    d_out = (out - y) * sigmoid_d(out)
    dW2 = ______         # A COMPLETER : h.T @ d_out / 4
    db2 = np.mean(d_out, axis=0, keepdims=True)

    d_h = d_out @ W2.T * sigmoid_d(h)
    dW1 = ______         # A COMPLETER : X.T @ d_h / 4
    db1 = np.mean(d_h, axis=0, keepdims=True)

    # Mise a jour
    W2 -= alpha * dW2; b2 -= alpha * db2
    W1 -= alpha * dW1; b1 -= alpha * db1

    if epoch % 1000 == 999:
        preds = (out > 0.5).astype(int).flatten()
        acc = np.mean(preds == y.flatten())
        print(f"Epoch {epoch+1} : "
              f"loss={loss:.4f}, acc={acc:.0%}")

print(f"\nPredictions : {out.flatten().round(3)}")
print(f"Attendu     : {y.flatten().astype(int)}")
\end{lstlisting}
\end{pythonbox}

% ======================================================================
\section*{Exercices Python autonomes --- Niveau avanc\'{e}}
% ======================================================================

% ----------------------------------------------------------------------
\exercice{\exoavance\quad --- R\'{e}seau from scratch pour salaire\_eleve \hfill\normalfont\small\textsf{[Python autonome]}}

\'{E}crivez un programme Python complet qui :

\begin{enumerate}[label=\alph*)]
    \item Charge le dataset salaire (30 d\'{e}veloppeurs, 2 features :
        exp\'{e}rience et nb\_langages) et la cible binaire (seuil 45\,k\texteuro).

    \item \textbf{Normalise} les features (soustraire la moyenne, diviser par l'\'{e}cart-type).

    \item Impl\'{e}mente un r\'{e}seau \textbf{2-4-1} from scratch avec :
    \begin{itemize}
        \item Forward pass (sigmo\"{\i}de partout)
        \item Backward pass (backpropagation)
        \item Mise \`{a} jour par descente de gradient ($\alpha = 0{,}5$, 2000 \'{e}poques)
    \end{itemize}

    \item Affiche la \textbf{loss et l'accuracy} toutes les 500 \'{e}poques.

    \item Trace \textbf{deux graphiques} c\^{o}te \`{a} c\^{o}te :
    \begin{itemize}
        \item \`{A} gauche : courbe de loss par \'{e}poque
        \item \`{A} droite : nuage de points + r\'{e}gions de d\'{e}cision du r\'{e}seau
            (utiliser \texttt{meshgrid} + \texttt{contourf})
    \end{itemize}

    \item Comparez l'accuracy finale avec le \textbf{perceptron} du Jour~3 (${\sim}87$--$90\%$).
\end{enumerate}

\begin{attentionbox}[Consignes]
    \begin{itemize}
        \item Utilisez \textbf{uniquement} NumPy et Matplotlib.
        \item La normalisation est \textbf{essentielle} pour que le r\'{e}seau converge.
        \item Initialisez les poids al\'{e}atoirement (\texttt{np.random.randn * 0.5}).
    \end{itemize}
\end{attentionbox}

% ----------------------------------------------------------------------
\exercice{\exoavance\quad --- NumPy vs scikit-learn MLPClassifier \hfill\normalfont\small\textsf{[Python autonome]}}

\'{E}crivez un programme qui compare deux impl\'{e}mentations :

\begin{enumerate}[label=\alph*)]
    \item \textbf{M\'{e}thode 1 --- From scratch :} votre r\'{e}seau 2-4-1 de l'exercice~11.

    \item \textbf{M\'{e}thode 2 --- scikit-learn :}
    \begin{verbatim}
from sklearn.neural_network import MLPClassifier
from sklearn.preprocessing import StandardScaler
    \end{verbatim}
    Utilisez \texttt{MLPClassifier(hidden\_layer\_sizes=(4,),
    max\_iter=2000, activation='logistic', random\_state=42)}.

    \item Affichez un \textbf{tableau comparatif} :
    \begin{center}
    \renewcommand{\arraystretch}{1.2}
    \begin{tabular}{l|c|c}
        \toprule
        \textbf{M\'{e}thode} & \textbf{Accuracy} & \textbf{Nb param\`{e}tres} \\
        \midrule
        From scratch (2-4-1) & ? & ? \\
        MLPClassifier (4,) & ? & ? \\
        Perceptron (Jour 3) & ? & ? \\
        \bottomrule
    \end{tabular}
    \end{center}

    \item Tracez les \textbf{surfaces de d\'{e}cision} des 3 mod\`{e}les c\^{o}te \`{a} c\^{o}te
        (3 sous-figures avec \texttt{contourf}).

    \item R\'{e}pondez en commentaire : le r\'{e}seau est-il meilleur que le perceptron ?
        La fronti\`{e}re est-elle lin\'{e}aire ou courbe ? Combien de param\`{e}tres le
        r\'{e}seau a-t-il vs le perceptron ?
\end{enumerate}

\begin{intuitionbox}[Indice]
    Le perceptron a 3 param\`{e}tres ($w_1, w_2, b$). Le r\'{e}seau 2-4-1 a
    $2 \times 4 + 4 + 4 \times 1 + 1 = 17$ param\`{e}tres. Plus de param\`{e}tres
    = plus de capacit\'{e} \`{a} mod\'{e}liser des fronti\`{e}res complexes, mais aussi
    plus de risque de surapprentissage.
\end{intuitionbox}

% ======================================================================
\vfill
\separateur

\begin{center}
    \small\sffamily
    \textbf{R\'{e}partition :}\quad
    4 exercices \starfull{ipssiVert} fondamentaux \quad
    3 exercices \starfull{ipssiOrange}\starfull{ipssiOrange} interm\'{e}diaires \quad
    3 guid\'{e}s Python \starfull{ipssiOrange}\starfull{ipssiOrange} \quad
    2 autonomes \starfull{ipssiRouge}\starfull{ipssiRouge}\starfull{ipssiRouge}
\end{center}

\begin{center}
    {\footnotesize\sffamily\textcolor{gray}{%
        \textcopyright\ 2025 Mohamed EL AFRIT --- IPSSI \enspace\textbullet\enspace
        \url{www.mohamedelafrit.com} \enspace\textbullet\enspace
        CC BY-NC-SA 4.0
    }}
\end{center}

\end{document}
